% ID: SEL-D0043-OB-E04
% ESTADO: validada. S044, cotejada con examen y criterios.
\textbf{1a)} La proporción del idioma C es
$1-1/3-1/2=1/6$. Las cuotas de 60 alumnos son
$\boxed{20,\ 30,\ 10}$ para A, B y C, respectivamente.

\textbf{1b)} Si 14 alumnos son un tercio de la muestra,
el total es 42. Corresponden $\boxed{21}$ a B y $\boxed{7}$ a C.

\textbf{2)} La esperanza de la media muestral es la media
poblacional, así que $\boxed{\mu=7}$. Si las muestras se
seleccionan sin reemplazo, usamos la corrección por población
finita, con $\sigma^2$ definida con denominador $N$:
\[
4=\operatorname{Var}(\bar X)
=\frac{\sigma^2}{3}\frac{5-3}{5-1}=\frac{\sigma^2}{6}.
\]
Por tanto $\boxed{\sigma^2=24}$. Si se entiende el muestreo
con reemplazo e independiente, la corrección desaparece
y $4=\sigma^2/3$, lo que da $\boxed{\sigma^2=12}$.
